\chapter{Message framing}

\section{Identifiers}
The instrument accepts exactly two System Exclusive identifiers:

\begin{center}
\begin{tabular}{ll}
\toprule
identifier & meaning \\
\midrule
\bytes{50} & Matsushita (Technics) \\
\bytes{7E} & Universal Non-Real Time \\
\bottomrule
\end{tabular}
\end{center}

A message carrying any other identifier is tracked to its \bytes{F7} and discarded. No
error is reported for it.

\begin{caution}{There is no device or unit number after the identifier}
Many instruments place a device id immediately after the manufacturer identifier. The
SX-WSA1R does not. The byte following \bytes{50} is already the message-family byte
described in chapter~\ref{ch:messageset}. A unit field does appear later in some families,
but it is not where the convention would put it.
\end{caution}

\section{Termination}
A message ends at \bytes{F7}. If any other status byte in the range
\bytes{80}--\bytes{F6} arrives while a message is open, the message is abandoned and its
contents discarded.

\section{Length limit}
A message may not exceed 254 bytes. A longer message is rejected and, inside a bulk-dump
session, reported as a reception error.
